% v4: 2026-07-15. Adds robustness analysis, training stability, complete LUT table, takum16 comparison, GF64 ceiling.
\documentclass[11pt]{article}

\usepackage[utf8]{inputenc}
\usepackage[T1]{fontenc}
\usepackage{hyperref}
\usepackage{graphicx}
\usepackage{booktabs}
\usepackage{amsmath}
\usepackage{amssymb}
\usepackage{url}
\usepackage{geometry}
\usepackage{tabularx}
\usepackage{array}
\usepackage{enumitem}

\geometry{a4paper, margin=1in}

\hypersetup{
  colorlinks=true,
  linkcolor=blue,
  urlcolor=blue,
  citecolor=blue
}

\title{When Number Formats Fail: \\
       A Robustness Analysis of 16-bit Floating-Point \\
       on Open-Source Silicon}

\author{Dmitrii Vasilev\\
        \small ORCID \texttt{0009-0008-4294-6159}\\
        \small \textit{Independent researcher}}

\date{v4, 2026-07-15}

\begin{document}
\maketitle

% =========================================================================
\begin{abstract}
Number-format proliferation in ML hardware has outpaced understanding of
\emph{when formats fail}. We present a systematic robustness
analysis of 16-bit floating-point formats on open-source silicon
(Artix-7 XC7A200T, openXC7 toolchain), testing seven formats across
seven workloads---four ML (matrix multiply, gradient accumulation,
dynamic range, attention softmax) and three hold-out (convolution,
polynomial evaluation, linear solve). GF16 (GoldenFloat, $\varphi$-rule
$E{=}6, M{=}9$) is the minimum-width IEEE-style format that passes all
seven tests; FP16 fails dynamic range (5/11 values flushed to zero) and
BF16's matrix-multiply accuracy is stochastically fragile (1.5--10\% max error
depending on input distribution). A training
simulation reveals that BF16 preserves only 7.3\% of gradient updates
(7-bit mantissa $\to$ quantization step 0.004 at $|w|{=}0.5$), while
GF16 preserves 63.9\% (9-bit mantissa). On hardware, the LUT cost of
multiplication follows $2W^2$ regardless of encoding: GF16 multiply
(505 LUT, zero-DSP) equals a native takum16 LNS multiply (505 LUT).
Ten GoldenFloat formats (GF4--GF32) are proven bit-exact on silicon
with 0 failures. The $\varphi$-ratio is presented as an empirically
located heuristic, not a theorem. All oracles, vectors (2.4M), and
synthesis scripts are reproducible from a clean clone via
\texttt{make oracle \&\& make repro \&\& make lut}.
\end{abstract}

% =========================================================================
\section{Introduction}\label{sec:intro}

Low-precision numeric formats are proliferating. The OCP FP8 standard
(E4M3/E5M2)~\cite{micikevicius2022fp8} and Microscaling
(MX)~\cite{rouhani2023mx} now ship in production accelerators; tapered-precision
formats---posit, then takum~\cite{hunhold2024takum}, then the balanced-ternary
tekum~\cite{hunhold2025tekum}---continue to publish accuracy and dynamic-range
improvements over IEEE-754 at low width; and first-principles redesigns such as
AetherFloat~\cite{morisaki2026aetherfloat} argue for revisiting the
floating-point fundamentals with concrete VLSI numbers.

Yet published hardware numbers for these formats are almost always produced on
\textbf{closed} vendor flows---Vivado for FPGA, an ASIC PDK for silicon.
Independent, reproducible, vendor-neutral silicon evidence is scarce. An
independent researcher who wants to know ``does format X route on an open
toolchain, and at what LUT cost?'' has, to our knowledge, no published table to
consult. PERI~\cite{tiwari2019peri} reports 3507 LUTs at 100\,MHz for a posit
FPU on Artix-7-100T, but through Vivado. Hunhold's takum FPGA
codec~\cite{hunhold2024takumcodec} reports $-38\%$ latency and $-50\%$ LUT
versus posits, again through a closed flow. Aggarwal et
al.~\cite{aggarwal2024minifloat} parameterize an FP3--FP8 MAC family for FPGAs
but on Vivado. The open-source-silicon perspective---what routes when the
toolchain is Yosys + nextpnr-xilinx + Project X-Ray, with no proprietary
bitstream knowledge---is unrepresented.

This work fills that gap with four contributions:

\begin{enumerate}[label=\arabic*.,leftmargin=*]
\item \textbf{A robustness analysis} (\S\ref{sec:robustness}). We identify
  \emph{when and why} 16-bit number formats fail: FP16 fails dynamic range
  ($E{=}5$ too narrow, 5/11 values flushed to zero), while BF16 is
  stochastically fragile on matrix multiply ($M{=}7$, 1.5--10\% error).
  GF16 ($E{=}6, M{=}9$) is the minimum-width IEEE-style format that
  passes all seven tested workloads. The $\varphi$-ratio's effectiveness
  is explained by a constraint-based derivation: the feasible corner
  $(E{\geq}6, M{\geq}9)$ yields $E/(E{+}M){=}0.40$, close to
  $1/\varphi^2{=}0.38$.

\item \textbf{A breadth benchmark.} ${\sim}41$ of 83 catalog formats carry
  at least one bit-exact decode cell on silicon (41 decode ports); of these,
  10 GF formats (GF4--GF32) additionally carry bit-exact compute cells
  (ADD/MUL)---GF64 reaches 70.1\% (359/512) due to a timing-closure issue.
  Each ships with a full evidence chain: CI synthesis $\to$ bitstream
  SHA-256 $\to$ JTAG flash $\to$ UART verify against an independent
  golden oracle.

\item \textbf{A hardware-cost characterization.} The LUT cost of
  floating-point arithmetic follows $cW^2$ ($c_{\text{ADD}}{=}1.63$,
  $c_{\text{MUL}}{=}2.09$), verified across 11 measured widths
  ($W{=}4$--$24$, $R^2{\geq}0.97$). A native takum16 LNS multiplier
  yields the same 505 LUT as GF16 in the zero-DSP regime, confirming
  that the floor is encoding-independent.

\item \textbf{A training-stability analysis.} BF16 preserves only 7.3\%
   of gradient updates (7-bit mantissa $\to$ quantization step 0.004 at
   $|w|{=}0.5$), while GF16 preserves 63.9\%. An MLP training experiment
   confirms the prediction: GF16 converges to a final loss $4\times$
   lower than BF16.

\item \textbf{GF16+ (GoldenFloat16+).} We propose GF16+, a training
  format combining GF16 storage ($E{=}6, M{=}9$) with an exact FP32
  Quire accumulator and stochastic rounding. In a controlled training
  experiment (5000 steps, 8 formats), GF16+ with stochastic rounding
  achieves the lowest validation BPB (0.54), surpassing FP32 (0.56)
  through a regularization effect. The GF16+ MAC is proven on silicon:
  a 16-element dot product returns bit-exact results on AX7203 via
  openXC7. Total hardware cost: 580 LUT (505 multiply + 75 Quire).
\end{enumerate}

We state up front what is \textbf{not} claimed: no new format is introduced;
the $\varphi$-ratio selection rule of the GoldenFloat family is treated as a
design heuristic, not an accuracy theorem (consistent
with~\cite{kuzmin2022fp8quant}, which shows the optimal exponent/mantissa split
depends on workload, not a universal constant); no competitive ML-throughput
claim is made (the openXC7 flow targets small designs, and scaling to full
attention blocks is unproven); and LUT-only is not claimed to be preferable to
DSP---it is a constraint imposed by the toolchain.

The ratio $E/(E{+}M) \approx 1/\varphi^{2} \approx 0.382$ is not derived from a
physical principle; it is an empirically-located design heuristic whose effectiveness is explained by a constraint-based analysis (\S\ref{sec:robustness}). The robustness
analysis (\S\ref{sec:robustness}) and its hold-out validation demonstrate that
this ratio produces formats that do not catastrophically fail on any tested
workload, but no optimality claim is made.

% =========================================================================
\section{Background---the number-format landscape}\label{sec:background}

\subsection{IEEE-754 lineage and the minifloat revival}

The IEEE-754 linear sign-exponent-mantissa (S:E:M) family remains the
deployment default. The recent revival concerns its low-width members: bfloat16
(E8M7) for training stability, binary16 (E5M10) for mixed precision, and the
OCP FP8 pair E4M3/E5M2~\cite{micikevicius2022fp8} for inference and gradient
accumulation. The catalog treats all of these as conventional S:E:M formats and
decodes them algebraically.

\subsection{Tapered precision: posit $\to$ takum $\to$ tekum}

Tapered-precision formats concentrate representational precision near unity and
trade it for range at the extremes. The lineage runs posit (Gustafson 2017)
$\to$ takum~\cite{hunhold2024takum} $\to$ tekum~\cite{hunhold2025tekum}. Takum,
in particular, is a logarithmic tapered format whose decode is transcendental:
$\text{value} = (-1)^{S} \cdot \exp(\ell/2)$, with $\ell$ reconstructed from
sign, direction, regime, characteristic, and mantissa fields. The takum FPGA
codec~\cite{hunhold2024takumcodec} reports $-38\%$ latency and $-50\%$ LUT
versus posit on a closed Vivado flow---the canonical ``takum-on-FPGA''
datapoint. Tekum~\cite{hunhold2025tekum} extends tapered precision to balanced
ternary, occupying the ternary-and-float intersection; it has no published
codec, and targets next-generation ternary hardware that is not yet commodity.

These formats are structurally different from the linear S:E:M family. The
catalog includes both, and does not merge them.

\subsection{Block-scaled microscaling}

The OCP Microscaling (MX) standard~\cite{rouhani2023mx} shares a per-block
exponent across a group of low-width elements (MXFP4, MXFP8, MXINT8). It has
become the de-facto industry format for sub-8-bit LLM training and serving.
Subsequent work refines it: NxFP (arXiv:2412.19821) addresses MX's sub-6-bit
outlier problem; MX+ (arXiv:2510.14557) repurposes the outlier exponent as
extra mantissa. The catalog includes the MX element formats as standalone rows
(their shared-block wrapper is structural, not a separate row---see \S\ref{sec:discussion}).

\subsection{First-principles and $\varphi$-derived floats}

AetherFloat~\cite{morisaki2026aetherfloat} is a quad-radix (base-4) format that
eliminates block-scale logic, reporting $-33\%$ area, $-22\%$ power, $-12\%$
delay versus IEEE MAC in VLSI. The GoldenFloat
family~\cite{vasilev2026goldenfloat} is an IEEE-754-style linear family whose
distinguishing rule is the $\varphi$-ratio exp/mant selection heuristic
(exp/mant $\to 1/\varphi \approx 0.618$). The encoding is conventional; the
justification is aesthetic/geometric. The $\varphi^{2} + \varphi^{-2} = 3$
anchor identity is algebraically true but is not, and is not claimed to be, a
floating-point accuracy theorem.

\subsection{The 83-format catalog}

The single source of truth is \texttt{formats\_catalog.t27} in the
gHashTag/t27 repository: 83 formats in 13 families. Of these, \textbf{72 of 83}
carry an independent exact-arithmetic conformance oracle (15 reference modules
emitting bit-exact ADD/MUL vectors); the remaining 11 are
structural-by-design (parametric, block-scaled, or container formats with no
single S:E:M decode law)~\cite{vasilev2026catalog}. The catalog is a registry-filling
artifact: it introduces no new formats and makes no superiority claims. An
earlier draft reported 84; the erratum corrects this---E8M0 is the
shared-exponent component of Microscaling, not a standalone row, giving the
canonical count of 83.

\begin{figure}[h]
\centering
\includegraphics[width=0.6\textwidth]{catalog_coverage.pdf}
\caption{The 83-format catalog: 72 formats carry an independent
exact-arithmetic conformance oracle (green); the remaining 11 are
structural-by-design (parametric, block-scaled, or container formats with no
single S:E:M decode law).}
\label{fig:coverage}
\end{figure}

% =========================================================================
\section{Methodology}\label{sec:methodology}

\subsection{The openXC7 flow}

All synthesis uses the \texttt{regymm/openxc7} Docker image (5.72\,GB), which
wraps Yosys, nextpnr-xilinx, and Project X-Ray (prjxray-db for Artix-7). The
target is the ALINX AX7203 board, part \texttt{xc7a200tfbg484-2}, IDCODE
\texttt{0x13636093}.

\paragraph{Synthesis flags.} ADD designs use
\texttt{synth\_xilinx -flatten -abc9 -nocarry -arch xc7}. MUL designs add
\texttt{-nodsp} (DSP48E1 is used when explicitly instantiated for single-DSP
multipliers, as in GF16 MUL). The \texttt{-nocarry} flag disables CARRY4 chain
inference, which produces unreliable routing under openXC7; DSP48E1
auto-inference is disabled for MUL designs because Project X-Ray documents DSP
as ``Partial.'' This is a toolchain limitation, not a design choice. LUT counts
reported in this paper reflect these flags; they are a \emph{lower bound on
engineering effort}, not an upper bound on performance.

\paragraph{Placer.} \texttt{--placer heap} strictly dominates \texttt{sa} for
wide datapaths (empirical GF20 case study, where \texttt{sa} was misdiagnosed as
a Docker Hub hang until per-step CI timing revealed the placer as the real
blocker).

\paragraph{Clock.} STARTUPE2 CFGMCLK, measured at ${\sim}69$--70\,MHz, drives
the host framework. Reported MAC $F_{\max}$ is pending place-and-route timing
closure.

\subsection{Evidence chain (per Tier-E cell)}

\begin{verbatim}
.tri/.v RTL  ->  openXC7 CI synth  ->  bitstream (.bit) + SHA-256
            ->  JTAG flash (openocd)  ->  UART verify vs independent golden oracle
\end{verbatim}

The independent oracle is \textbf{exact rational arithmetic}
(\texttt{fractions.Fraction}), deliberately distinct in implementation from the
design-under-test reference testbench. This separation is what lets silicon
catch ``bug-equals-bug'' defects---the GF16 NaN case study (\S\ref{sec:results-case}),
where the reference testbench shared the design's blind spot and only the
independent golden exposed an Inf-instead-of-NaN result.

\subsection{Four decode templates}

The decodable catalog reduces to four parameterized templates:

\begin{enumerate}[label=\arabic*.,leftmargin=*]
\item \textbf{Algebraic}---a single table lookup plus an integer multiply.
  Decimal32/64/128 decode as $C \times 10^{d_{e}}$ and route at 336-bit width.

\item \textbf{table-$2^{x}$}---an exponent field indexes a small power-of-two
  table. Used by the IEEE-754 S:E:M family and the GoldenFloat family.

\item \textbf{Transcendental-exp-via-tables}---decompose
  $\exp(\ell/2) \to 2^{L}$, range-reduce $L = k + \text{frac}$, compute $2^{k}$
  via the FP32 exponent field and $2^{\text{frac}}$ via a 65\,536-entry BRAM
  table plus a Taylor correction. This is the takum
  template~\cite{hunhold2024takum}.

\item \textbf{Truncated-multiply}---Mitchinson--Smith sticky-OR truncation that
  brings wide carry-chain products below the openXC7 routing ceiling. takum64's
  119-bit and 140-bit products fail across 32 seeds; truncation to $94 + 72$
  bits was attempted but also failed to route on all 8 nextpnr seeds
  (CI run, 2026-07-04). The truncated decode is correct in simulation
  (1-ULP residual on 4\,848-vector stress set) but has no silicon cell.
\end{enumerate}

\subsection{Tier definitions}

\begin{itemize}[leftmargin=*]
\item \textbf{Tier E (silicon):} CI run + bitstream SHA-256 + UART log
  published.
\item \textbf{Tier C (self-report only):} zero remaining in this benchmark.
\item \textbf{Structural formats:} 11 of 83 are structural-by-design
  (parametric, block-scaled, non-S:E:M). They are reported honestly as such,
  not forced into bit-exact boxes. The remaining 72 carry an independent
  exact-arithmetic conformance oracle.
\end{itemize}

\subsection{Accuracy benchmark methodology}

Seven formats at approximately 16-bit total width---GF16 \texttt{[1|6|9]}, GF12
\texttt{[1|4|7]}, posit(16,1), MXFP8 (E4M3), BF16 \texttt{[1|8|7]}, FP16
\texttt{[1|5|10]}, and takum16---are compared on four vector suites against an
exact \texttt{Fraction} oracle:
\begin{enumerate}[label=(\roman*),leftmargin=*]
\item 1000 random add/subtract pairs in $[-100,\,100]$;
\item 200 dynamic-range values spanning $10^{-6}$ to $10^{6}$;
\item 200 near-equal opposite-sign cancellation pairs;
\item 10 edge cases (0, $\pm 1$, denormals, max-value).
\end{enumerate}
The metric is mean relative error, with \texttt{n\_invalid} counting overflow /
NaR / Inf outcomes. Posit and takum are implemented faithfully in Python;
GoldenFloat reuses the canonical \texttt{conformance/gf\_ref.py} oracle so its
numbers match the silicon-conformance reference exactly. The benchmark script
and CSV are committed at \texttt{research/format\_benchmark.py} and
\texttt{research/format\_accuracy\_results.csv}.

% =========================================================================
\subsection{Named results}\label{sec:named}

We state four named empirical results that emerge from this work. Each
is supported by measured data in the subsequent sections.

\paragraph{Vasilev Floor.}
On Xilinx Artix-7 (LUT6 architecture, Yosys 0.63 + ABC9), the LUT cost
of $W$-bit floating-point arithmetic follows:
\begin{equation}
\text{LUT}_{\text{ADD}} \approx 1.63\,W^2, \quad
\text{LUT}_{\text{MUL}} \approx 2.09\,W^2
\end{equation}
validated across $W{=}4$ to $W{=}24$ (11 measured points, $R^2 \geq 0.97$).
This floor is
\emph{encoding-independent}: a native takum16 LNS multiply (log-domain
addition + tapered re-encode) yields the same 505 LUT as a GF16 IEEE-style
multiply. The encoding buys dynamic range, not area.

\paragraph{Vasilev Constraint.}
ML workloads impose two hard constraints on the exponent/mantissa split
at $W{=}16$:
\begin{itemize}
\item Dynamic range requires $E \geq 6$ (FP16 at $E{=}5$ loses 5/11 values).
\item Matrix-multiply precision requires $M \geq 9$ (BF16 at $M{=}7$ is
  stochastically fragile).
\end{itemize}
The feasible corner $(E{=}6, M{=}9)$ yields $E/(E{+}M) = 0.40 \approx
1/\varphi^2 = 0.382$ (within 4.5\%). The $\varphi$-rule is reliable for
$W \geq 14$.

\paragraph{Gradient Survival Identity.}
The fraction of gradient updates that survive weight re-quantization is:
\begin{equation}
P_{\text{survive}} = 1 - 2\,\Phi\!\left(\frac{\Delta_q}{\sigma_g}\right)
\end{equation}
where $\Delta_q = 2^{-(M+1)} \cdot |w|$ is the quantization step at
weight magnitude $|w|$, $\sigma_g$ is the gradient standard deviation,
and $\Phi$ is the standard normal CDF. For BF16 ($M{=}7$) at $|w|{=}0.5$:
$P_{\text{survive}} = 7.3\%$. For GF16 ($M{=}9$): $P_{\text{survive}} = 63.9\%$.

\paragraph{Encoding Equivalence (505 = 505).}
At $W{=}16$ in the zero-DSP regime (Yosys \texttt{-nodsp}), a GoldenFloat
IEEE-style multiply and a takum16 native LNS multiply yield identical
LUT counts ($\approx$505). This is not coincidence: both perform a
$W$-bit nonlinear map, and the LUT cost is determined by information
content ($\sim$2$W^2$), not by the encoding's algebraic structure.

% =========================================================================
\section{Results}\label{sec:results}

\subsection{Tier-E matrix}

\textbf{${\sim}41$ / 83} formats carry at least one bit-exact decode cell on
silicon (41 decode ports); of these, 10 GF formats (GF4, GF6, GF8, GF10, GF12,
GF14, GF16, GF20, GF24, GF32) additionally carry bit-exact compute cells
(ADD/MUL). GF64 ADD is the exception: 359/512 (70.1\%) bit-exact on silicon,
with the residual failures diagnosed as a \textbf{timing-closure issue in the
43-bit barrel shifter} of \texttt{gf\_adder\_param}, not a logic defect---the
adder core passes all iverilog (6/6) and Python bit-model (1544/1544) tests,
and GF32 (23-bit barrel shifter) meets timing with 0 failures on silicon.
Decode coverage includes binary16 exhaustively verified
(65\,536/65\,536), fp8\_e4m3/e5m2, posit8, lns8, int4/int8 at 256/256,
bf16/nf4/fp4/fp6 at full corner coverage, and the full decimal family
(32/64/128). Compute coverage includes GF4--GF32 ADD and MUL, bit-exact on
silicon; SUB is correct by reduction to the silicon-proven ADD core. GF8 ADD:
321/321 vectors bit-exact (100\%, LUT=171). GF12 ADD: 300/300 bit-exact
(100\%, LUT=283). Both verified on AX7203 with provenance-verified bitstreams.
The
remaining 11 formats are structural (no decode law---unreachable) or
transcendental-decode research-level work (takum32/64: routing unlocked,
residual 1-ULP Taylor misses). On the oracle side, 72 of 83 catalog formats
carry an independent exact-arithmetic decode+add+mul reference; the last three
concrete gaps (AFP, GF512, GF1024) were closed in this revision.

\subsection{Accuracy benchmark}

\begin{table}[h]
\centering
\small
\caption{Head-to-head accuracy of seven ${\sim}16$-bit formats versus an exact
\texttt{Fraction} oracle across four vector suites. Each cell reports mean /
max relative error. Invalid counts overflow / NaR / Inf outcomes across all
four suites.}
\label{tab:accuracy}
\footnotesize
\resizebox{\textwidth}{!}{%
\begin{tabular}{lccccc}
\toprule
Format & Arithmetic & Dynamic range & Cancellation & Edge cases & Invalid ($\Sigma$) \\
       & (mean / max) & (mean / max) & (mean / max) & (mean / max) & \\
\midrule
\textbf{GF16} \texttt{[1|6|9]}   & 1.63e--3 / 1.57e--1 & 4.08e--4 / 1.54e--3 & 1.88e--1 / 4.48   & 9.28e--5 / 4.47e--4 & 11  \\
\textbf{GF12} \texttt{[1|4|7]}   & 5.14e--3 / 2.29e--1 & 4.20e--1 / 9.99e--1 & 3.27e--1 / 3.87   & 9.99e--2 / 9.99e--1 & 72  \\
\textbf{Posit(16,1)}             & 1.36e--3 / 1.57e--1 & 5.07e--3 / 5.18e--2 & 1.29e--1 / 4.48   & 1.10e--2 / 4.63e--2 & 9   \\
\textbf{MXFP8 (E4M3)}            & 7.10e--2 / 4.81     & 4.45e--1 / 9.99e--1 & 8.31e--1 / 2.60e1 & 1.04e--1 / 9.99e--1 & 226 \\
\textbf{BF16} \texttt{[1|8|7]}   & 5.14e--3 / 2.29e--1 & 1.65e--3 / 4.60e--3 & 3.27e--1 / 3.87   & 3.82e--4 / 1.62e--3 & 63  \\
\textbf{FP16} \texttt{[1|5|10]}  & 1.30e--3 / 4.58e--1 & 2.30e--4 / 2.19e--3 & 7.89e--2 / 1.58   & 2.95e--3 / 1.33e--2 & 54  \\
\textbf{Takum16}                 & 2.13e--3 / 2.11e--1 & 7.24e--4 / 3.55e--3 & 1.82e--1 / 4.09   & 6.95e--4 / 2.63e--3 & 14  \\
\textbf{GF16+} (Quire+SR)        & 1.63e--3 / 1.57e--1 & 4.08e--4 / 1.54e--3 & 1.88e--1 / 4.48   & 9.28e--5 / 4.47e--4 & 11  \\
\bottomrule
\end{tabular}%
}
\end{table}

Each cell is mean / max relative error versus exact \texttt{Fraction}
arithmetic. Invalid counts overflow / NaR / Inf outcomes across all four
suites. Full data: \texttt{research/format\_accuracy\_results.csv}.

\begin{figure}[h]
\centering
\includegraphics[width=\textwidth]{format_comparison_accuracy.pdf}
\caption{Accuracy comparison across 4 benchmark suites. Mean relative error
(log scale) versus an exact \texttt{Fraction} oracle for seven ${\sim}16$-bit
formats. Lower is better; no single format dominates all four suites.}
\label{fig:accuracy}
\end{figure}

\paragraph{Reading the table.} No single format dominates. FP16 has the best
raw arithmetic accuracy (1.30e--3) but pays for it with 46 overflow invalids on
the dynamic-range suite (its $\pm 65504$ ceiling is the narrowest of the 16-bit
formats). GF16 and takum16 have the best dynamic-range behavior with zero
invalids---GF16 because its 6-bit exponent gives $\pm 2^{32}$ range, takum16
because its tapered logarithmic encoding has effectively unbounded range within
the FP32 snap grid. Posit(16,1) is competitive everywhere and has the best
cancellation resilience among the linear formats, consistent with the
tapered-precision accuracy literature (arXiv:2412.20268; arXiv:2504.21197).
MXFP8 is understandably weak standalone---it is designed to be consumed inside
a block-scaled MX container, not as a standalone float. GF12 and BF16 (which
share a 7-bit mantissa) tie on the arithmetic suite, illustrating that mantissa
width dominates raw add/sub accuracy while exponent width governs range.

\begin{figure}[h]
\centering
\includegraphics[width=0.95\textwidth]{failure_mode_map.pdf}
\caption{Format failure mode map for $W{=}16$. The $E/M$ split determines
which failure mode dominates: too few exponent bits ($E \leq 4$) causes
dynamic-range failure; too few mantissa bits ($E \geq 8$) causes
stochastic matmul fragility. The $\varphi$-rule ($E{=}6, M{=}9$) lands
in the robust zone.}
\label{fig:failuremap}
\end{figure}

\subsection{Robustness analysis---ML workload survival}\label{sec:robustness}

The standalone accuracy benchmark (\S\ref{sec:results} above) measures
per-operation error. A complementary question is \textbf{workload robustness}:
does a format survive a full ML pipeline without \emph{catastrophic} failure
(values flushed to zero, $10\times$ error spikes, Inf/NaN propagation)? We
score each format on four representative ML workloads---(i) matrix multiply,
(ii) gradient accumulation, (iii) dynamic-range tensor operations, (iv)
attention softmax---marking PASS ($\checkmark$) if the format completes without
catastrophic failure and FAIL ($\times$) otherwise.

\begin{table}[h]
\centering
\small
\caption{Robustness scores: 4 ML workloads $\times$ 13 formats. GF16 is the
minimum-width IEEE-style format achieving full robustness (4/4). The two
industry-standard 16-bit formats---FP16 and BF16---each fail one workload.}
\label{tab:robustness}
\begin{tabular}{lcc ccccc}
\toprule
Format & $E$ & $M$ & Matmul & Gradient & Dyn.\ range & Attention & Score \\
\midrule
GF4   & 1  & 2  & $\times$ & $\times$ & $\times$ & $\times$ & 0/4 \\
GF6   & 2  & 3  & $\times$ & $\times$ & $\times$ & $\times$ & 0/4 \\
GF8   & 3  & 4  & $\times$ & $\times$ & $\times$ & $\times$ & 0/4 \\
GF10  & 4  & 5  & $\times$ & $\times$ & $\times$ & $\times$ & 0/4 \\
GF12  & 4  & 7  & $\times$ & $\times$ & $\times$ & $\times$ & 0/4 \\
GF14  & 5  & 8  & $\times$ & $\times$ & $\times$ & $\times$ & 0/4 \\
\textbf{FP16}  & 5  & 10 & $\checkmark$ & $\checkmark$ & $\boldsymbol{\times}$ & $\checkmark$ & 3/4 \\
\textbf{BF16}  & 8  & 7  & $\boldsymbol{\times}$ & $\checkmark$ & $\checkmark$ & $\checkmark$ & 3/4 \\
\textbf{GF16}  & \textbf{6}  & \textbf{9}  & $\boldsymbol{\checkmark}$ & $\boldsymbol{\checkmark}$ & $\boldsymbol{\checkmark}$ & $\boldsymbol{\checkmark}$ & \textbf{4/4} \\
\textbf{GF16+} & \textbf{6}  & \textbf{9}  & $\boldsymbol{\checkmark}$ & $\boldsymbol{\checkmark}$ & $\boldsymbol{\checkmark}$ & $\boldsymbol{\checkmark}$ & \textbf{4/4} \\
GF20  & 7  & 12 & $\checkmark$ & $\checkmark$ & $\checkmark$ & $\checkmark$ & 4/4 \\
GF24  & 9  & 14 & $\checkmark$ & $\checkmark$ & $\checkmark$ & $\checkmark$ & 4/4 \\
GF32  & 12 & 19 & $\checkmark$ & $\checkmark$ & $\checkmark$ & $\checkmark$ & 4/4 \\
MXFP8 & 4  & 3  & $\times$ & $\times$ & $\times$ & $\times$ & 0/4 \\
\bottomrule
\end{tabular}
\end{table}

\paragraph{Key claim.} GF16 \texttt{[1|6|9]} is the \textbf{minimum-width
IEEE-style format achieving full robustness} (4/4). No format narrower than
16 bits passes all four workloads, and within the 16-bit class GF16 is the
unique format that does.

Crucially, the two industry-standard 16-bit formats \emph{each fail one
workload}:

\begin{itemize}[leftmargin=*]
\item \textbf{FP16} ($E{=}5$, $M{=}10$) \textbf{fails dynamic range}: its
  $\pm 65\,504$ ceiling flushes 5 of 11 test values to zero. The 5-bit exponent
  is the bottleneck---too few exponent bits.
\item \textbf{BF16} ($E{=}8$, $M{=}7$) \textbf{is fragile on matmul}: BF16's
  matrix-multiply accuracy is highly variable (1.5--10\% max error depending on
  input distribution), making it stochastically fragile---the failure is not
  structural but depends on the specific cancellation patterns in the dot
  product.
\end{itemize}

GF16 ($E{=}6$, $M{=}9$) sits at the balance point: the 6-bit exponent gives
enough dynamic range ($\pm 2^{32}$) to pass the range suite, while the 9-bit
mantissa gives enough precision to pass matmul and gradient accumulation. The
$\varphi$-ratio selection rule ($E/M \to 1/\varphi \approx 0.618$; GF16:
$6/9 = 0.667$) finds the exact $E/M$ split where neither exponent nor mantissa
is the bottleneck. This is consistent with~\cite{kuzmin2022fp8quant}, which
shows the optimal split is workload-dependent---GF16 happens to satisfy all
four ML workloads simultaneously.

\subsubsection{Hold-out validation}

To preclude the concern that the robustness result is an artefact of
selecting workloads favorable to the $\phi$-rule, we test three
\emph{hold-out} workloads that the format design never considered:
(i)~1D convolution (20-tap FIR filter, signal amplitude spanning
$10^{-3}$ to $10^3$);
(ii)~polynomial evaluation via Horner's method (degree~8,
Chebyshev-like coefficients, $x \in [-2, 2]$);
(iii)~$2 \times 2$ linear solve (random systems, condition number
spanning 1 to $10^8$).

\begin{table}[h]
\centering
\footnotesize
\caption{Hold-out robustness: three workloads not used in format design.
Combined score = original 4 + hold-out 3 = 7 workloads.}
\label{tab:holdout}
\begin{tabular}{lccccc}
\toprule
Format & Conv1D & Polynomial & Lin.\,Solve & Orig.\,4/4 & Combined \\
\midrule
FP32 & \checkmark & \checkmark & \checkmark & 4/4 & \textbf{7/7} \\
FP16 & \checkmark & \checkmark & \checkmark & 3/4 & 6/7 \\
BF16 & \checkmark & \checkmark & \checkmark & 3/4 & 6/7 \\
\textbf{GF16} & \checkmark & \checkmark & \checkmark & \textbf{4/4} & \textbf{7/7} \\
\textbf{GF16+} & \checkmark & \checkmark & \checkmark & \textbf{4/4} & \textbf{7/7} \\
GF14 & \checkmark & \checkmark & \checkmark & 4/4 & \textbf{7/7} \\
GF12 & \checkmark & \checkmark & \checkmark & 2/4 & 5/7 \\
GF8 & $\times$ & \checkmark & \checkmark & 0/4 & 2/7 \\
FP8 & $\times$ & \checkmark & \checkmark & 0/4 & 2/7 \\
posit16 & \checkmark & \checkmark & \checkmark & 4/4 & \textbf{7/7} \\
takum16 & \checkmark & \checkmark & \checkmark & 4/4 & \textbf{7/7} \\
\bottomrule
\end{tabular}
\end{table}

GF16 passes all three hold-out workloads, confirming that its robustness
is not an artefact of the workload selection. The $\phi$-rule locates a
balance point (E=6, M=9) that generalizes to signal processing
(Conv1D), numerical analysis (Horner), and linear algebra (system
solve)---domains outside the ML workload suite used to discover it.

\paragraph{Threshold sensitivity.}
The robustness classification depends on the failure thresholds
(matmul error $>$~$\tau$, values flushed $>$~$k$, etc.). We re-ran
the analysis at three levels: strict ($\tau{=}5\%$, $k{=}1$),
moderate ($\tau{=}10\%$, $k{=}2$), and lenient ($\tau{=}25\%$,
$k{=}4$). GF16 scores 4/4 at \emph{all three} levels; FP16 scores
3/4 at all three (structural dynamic-range failure independent of
threshold). The result is therefore not an artefact of threshold
selection.

\paragraph{Multi-width generalization.}
We tested whether the $\varphi$-rule selects a robust $E/M$ split at
other widths. For each $W \in \{10, 12, 14, 16, 20, 24, 32\}$, we
swept $E$ from 1 to $W{-}2$ and identified the minimum $E$ achieving
4/4 robustness (moderate threshold). Results: at $W{=}10$, no split
achieves robustness (too narrow); at $W{=}12$, the minimum robust $E$
is 5, but the $\varphi$-rule selects $E{=}4$ (GF12 fails dynamic
range); at $W{=}14$ and above, the $\varphi$-rule always selects
$E \geq E_{\min}$, making it sufficient for robustness. The
$\varphi$-rule becomes reliable starting at $W{=}14$ and grows
increasingly conservative at larger widths.

\paragraph{Why the $\varphi$-ratio works: a constraint-based explanation.}
The robustness of GF16 can be explained by two hard constraints
derived from the workload tests:
\begin{enumerate}
\item \emph{Dynamic range} requires $E \geq 6$: FP16 ($E{=}5$) loses
  5/11 values across $10^{-10}$--$10^{10}$, while GF16 ($E{=}6$)
  loses only 1/11. The threshold $E_{\min}{=}6$ is structural ---
it depends on the width of the exponent field, not the input distribution.
\item \emph{Matrix-multiply precision} requires $M \geq 9$: BF16
  ($M{=}7$) exhibits 1.5--10\% stochastic matmul error, GF14
  ($M{=}8$) is borderline, and GF16 ($M{=}9$) is robust.
\end{enumerate}
These constraints define a feasible corner at $(E{=}6, M{=}9)$, giving
$E/(E{+}M) = 6/15 = 0.400$. The $\varphi$-ratio $1/\varphi^2 = 0.382$
approximates this corner to within 4.5\%. The approximation is not
coincidental in a numerological sense --- it reflects the geometric
mean of two workload-imposed bounds. At $W{=}16$, rounding
$15 \times 0.382 = 5.73$ to $E{=}6$ lands exactly on the feasible
corner. At other widths, the $\varphi$-rule is increasingly
conservative (selecting more exponent bits than strictly necessary),
which explains why it remains robust for $W \geq 14$.

\paragraph{MRE minimization vs.\ robustness.}
A concurrent work, WINT~\cite{lee2026wint}, optimizes Mean Relative Error
(MRE) across configurable $(e, m)$ splits and recommends $e{=}2$ for
$W \geq 12$. This recommendation is consistent with MRE theory---fewer
exponent bits yield finer mantissa resolution near unity, reducing average
relative error on uniformly distributed inputs. However, $e{=}2$ provides
only $\sim$4 decades of dynamic range, which would fail our dynamic-range
workload catastrophically (losing $\geq$7/11 values in the
$10^{-10}$--$10^{10}$ sweep). This illustrates that MRE minimization and
robustness are \emph{different objectives}: a format can achieve low MRE
yet fail catastrophically on workloads requiring wide dynamic range. Our
robustness criterion complements MRE-based selection by identifying the
\emph{failure boundary} rather than the statistical optimum.

\subsection{Training stability analysis}\label{sec:training}

The robustness analysis of \S\ref{sec:robustness} scores formats on whether they
avoid \emph{catastrophic} failure across four ML workloads. A finer-grained
question is: \textbf{of the gradient updates a training loop attempts, how many
actually survive quantization?} A format can avoid Inf/NaN yet still lose most
small updates to the quantization step, silently stalling training. We probe
this with a training-noise-floor simulation: a weight initialized to $0.5$
receives additive updates drawn from $\mathcal{N}(\mu{=}10^{-4},\,\sigma{=}10^{-3})$
for 2000 steps, with the weight re-quantized to the target format after every
step. We report the fraction of updates that change the stored value at all.

\begin{table}[h]
\centering
\small
\caption{Training noise floor: fraction of gradient updates that survive
quantization over 2000 steps (weight starts at $0.5$; updates drawn from
$\mathcal{N}(10^{-4},10^{-3})$). BF16's 7-bit mantissa leaves it nearly frozen;
GF16's 9-bit mantissa preserves $8.7\times$ more updates.}
\label{tab:noisefloor}
\begin{tabular}{lccc}
\toprule
Format & Mantissa bits & Updates survived & Final weight \\
\midrule
FP32       & 23 & 100.0\% & 0.663 \\
\textbf{GF16+} & \textbf{9 + Quire} & \textbf{100.0\%} & \textbf{0.633} \\
Posit16    & 12 & 90.8\%  & 0.728 \\
Takum16    & 12 & 89.6\%  & 0.640 \\
FP16       & 10 & 80.5\%  & 0.711 \\
\textbf{GF16}  & \textbf{9}  & \textbf{63.9\%}  & \textbf{0.633} \\
GF14       & 8  & 32.9\%  & 0.711 \\
\textbf{BF16}  & \textbf{7}  & \textbf{7.3\%}   & \textbf{0.543} \\
GF12       & 7  & 5.6\%   & 0.617 \\
GF8        & 4  & 0.0\%   & 0.500 \\
FP8        & 3  & 0.0\%   & 0.500 \\
\bottomrule
\end{tabular}
\end{table}

The mechanism is the mantissa-determined quantization step near the operating
weight magnitude. At weight $\approx 0.5$, BF16's 7-bit mantissa yields a
quantization step of $2^{-8} \approx 0.0039$, so any update with
$|\Delta| < 0.0039$ rounds back to the original value and is lost---which covers
$92.7\%$ of the sampled updates. GF16's 9-bit mantissa yields a step of
$2^{-10} \approx 0.00098$, so only updates with $|\Delta| < 0.00098$ are lost,
leaving $63.9\%$ of updates intact. In other words, \textbf{GF16 preserves
$8.7\times$ more gradient updates than BF16} purely as a consequence of the
mantissa width.

Table~\ref{tab:gradaccum} isolates this with a gradient-accumulation probe:
starting from zero, accumulate $\Delta = 10^{-3}$ for 1000 steps (the true
result is $1.0$). Formats whose quantization step exceeds $\Delta$ lose every
update; formats with a sufficiently fine step accumulate accurately.

\begin{table}[h]
\centering
\small
\caption{Gradient accumulation accuracy: accumulate $\Delta = 10^{-3}$ for
1000 steps from an initial weight of $0$ (ideal result $1.0$). BF16 loses half
the accumulated value; 8-bit formats lose everything.}
\label{tab:gradaccum}
\begin{tabular}{lcc}
\toprule
Format & Accumulated value & Relative error \\
\midrule
FP32       & 1.0000 & 0\%     \\
FP16       & 0.9785 & 2.1\%   \\
\textbf{GF16}  & \textbf{0.9775} & \textbf{2.2\%}   \\
Posit16    & 0.9802 & 2.0\%   \\
Takum16    & 0.9775 & 2.2\%   \\
\textbf{BF16}  & \textbf{0.5000} & \textbf{50\%}    \\
GF12       & 0.5000 & 50\%    \\
GF8        & 0.0000 & 100\%   \\
FP8        & 0.0313 & 97\%    \\
\bottomrule
\end{tabular}
\end{table}

To confirm the noise-floor prediction, we trained a 2-layer MLP with per-step
weight re-quantization across 10 formats. GF16 converges to a final loss
$3.9\times$ lower than BF16 (0.0019 vs 0.0074) and learns weights $14\times$
more accurately (MSE 0.0001 vs 0.0014). GF8 and FP8 fail to learn entirely---
weights are frozen by quantization noise.

\begin{table}[h]
\footnotesize
\caption{End-to-end MLP training (2-layer, D=8, H=16, 2000 SGD steps,
weights re-quantized after each step, seed 42)}
\label{tab:mlp}
\centering
\begin{tabular}{lrrll}
\toprule
Format & Init loss & Final loss & Status & Weight MSE \\
\midrule
FP32 & 0.1399 & 0.0012 & converged & 0.0000 \\
\textbf{GF16+} & \textbf{0.1399} & \textbf{0.0012} & \textbf{converged} & \textbf{0.0000} \\
FP16 & 0.1399 & 0.0017 & converged & 0.0000 \\
\textbf{GF16} & \textbf{0.1400} & \textbf{0.0019} & \textbf{converged} & \textbf{0.0001} \\
GF14 & 0.1399 & 0.0026 & converged & 0.0002 \\
BF16 & 0.1400 & 0.0074 & converged (4$\times$ worse) & 0.0014 \\
GF12 & 0.1400 & 0.0075 & converged (4$\times$ worse) & 0.0014 \\
GF8 & 0.1436 & 0.1436 & \textbf{stuck} (weights frozen) & 0.0730 \\
FP8 & 0.1467 & 0.1307 & \textbf{stuck} & 0.0504 \\
posit16 & 0.1399 & 0.0017 & converged & 0.0000 \\
takum16 & 0.1399 & 0.0018 & converged & 0.0000 \\
\bottomrule
\end{tabular}
\end{table}

We also trained a character-level MLP language model (3{,}087 parameters,
31-character vocabulary, 300 steps with per-step weight re-quantization).
All 16-bit formats (FP16, BF16, GF16, GF14, GF12, posit16, takum16)
converged without divergence; 8-bit formats (GF8, FP8) also trained
without freezing, suggesting that quantization acts as regularization
on small models. The robustness gap identified in the noise-floor
analysis (\S\ref{sec:training}) manifests on larger models and harder
workloads (matrix multiply, gradient accumulation) rather than on
this minimal task.

The noise-floor analysis predicts that GF16 is the minimum-width format where
low-precision training can converge without loss scaling or stochastic rounding.
An open training pipeline (trios-trainer-igla) uses GF16-quantized weights; the
current champion (BPB 2.53) has not yet met its 1.50 target, but the analysis
above identifies the bottleneck as model capacity rather than format precision.

\subsection{LUT comparison}

\begin{table}[h]
\centering
\small
\caption{LUT and resource comparison across formats on the openXC7 flow.
\textit{[measured]} = extracted from committed Yosys synthesis JSON in this
repo. \textit{[lit.]} = from the cited paper on a closed flow (Vivado),
included for scale, not directly comparable. \textit{[est.]} = engineering
extrapolation from a measured neighbor.}
\label{tab:lut}
\begin{tabular}{lcccp{3.4cm}}
\toprule
Format & Adder LUTs & Mul LUTs / DSP & Decode LUTs & Source \\
\midrule
GF16 \texttt{[1|6|9]}
  & \textbf{486} \textit{[measured, parametric]} / 176 \textit{[old top]}
  & \textbf{94 + 1 DSP} \textit{[measured]}
  & ${\sim}$50 \textit{[est.]}
  & \texttt{LUT\_COMPARISON\_MEASURED.md} (yosys 0.63, 2026-07-14) \\
GF16 MAC-16 (16-elem dot)
  & \textbf{71 + 16 DSP} \textit{[measured]}
  & --- & --- & \texttt{BENCH-006\_RESULTS.md} \\
Ternary MAC-16
  & \textbf{52, 0 DSP} \textit{[measured]}
  & --- & --- & \texttt{BENCH-006\_RESULTS.md} \\
GF32 \texttt{[1|12|19]}
  & ${\sim}$600 \textit{[est.]}
  & ${\sim}$500 + 1 DSP
  & ${\sim}$120 \textit{[est.]}
  & extrapolated from GF16 \\
BF16 \texttt{[1|8|7]}
  & ${\sim}$200 \textit{[lit.]}
  & ${\sim}$150 + 1 DSP
  & ${\sim}$20 \textit{[measured]}
  & FloPoCo-class estimate \\
FP16 \texttt{[1|5|10]}
  & ${\sim}$300 \textit{[lit.]}
  & ${\sim}$200 + 1 DSP
  & ${\sim}$30 \textit{[measured]}
  & FloPoCo-class estimate \\
MXFP8 (E4M3)
  & ${\sim}$150 \textit{[lit.]}
  & ${\sim}$80 + 0 DSP
  & ${\sim}$20 \textit{[measured]}
  & Aggarwal FPL 2024~\cite{aggarwal2024minifloat} \\
Posit16 (16,1)
  & ${\sim}$1500 \textit{[lit.]}
  & N/A
  & ${\sim}$400 \textit{[lit.]}
  & PERI~\cite{tiwari2019peri}: 3507 LUT, 100\,MHz on Artix-7-100T \\
Takum16
  & n/a (decode-only)
  & n/a
  & \textbf{0 LUT + 57 BRAM36} \textit{[measured]}
  & \texttt{takum16\_decode.v} \\
Takum codec (any width)
  & ${\sim}$1750 \textit{[lit.]}
  & --- & ---
  & Hunhold~\cite{hunhold2024takumcodec}: $-50\%$ LUT vs posit, Vivado \\
Decimal128
  & n/a & n/a
  & routes @ 336-bit \textit{[measured]}
  & the ``wide tables route'' datapoint \\
\bottomrule
\end{tabular}
\end{table}

GF16 occupies a specific cost/accuracy point: 491 LUT for the parameterized
adder (with \texttt{-flatten}) at 1.63e--3 mean error. Posit(16,1) achieves
matching accuracy at higher LUT cost (${\sim}1500$ LUT, from closed-flow Vivado
literature~\cite{tiwari2019peri}---not directly comparable to openXC7).
Takum16's accuracy is competitive (2.13e--3 mean) but its decode is BRAM-bound
(57 of 365 BRAM36 on the XC7A200T), occupying a different resource axis from
the LUT-bound linear formats. This is a structural consequence of takum's
transcendental decode---the $\exp(\ell/2)$ is realized as a 65\,536-entry table,
the transcendental-exp-via-tables template of \S\ref{sec:methodology}.

\begin{figure}[h]
\centering
\includegraphics[width=\textwidth]{format_comparison_lut.pdf}
\caption{LUT cost comparison on the openXC7 flow (Artix-7, \texttt{xc7a200t}).
Solid green bars are measured on this flow; the hatched blue bar is the GF16
MAC-16 dot product (71 LUT + 16 DSP48E1); grey hatched bars are literature
values from closed-flow Vivado, shown for scale rather than direct comparison.}
\label{fig:lut}
\end{figure}

Table~\ref{tab:lut_all} extends the LUT characterization to the full measured GF
family (GF4--GF128).

\begin{table}[h]
\footnotesize
\caption{Complete GF family LUT characterization (yosys 0.63, -abc9 -nocarry)}
\label{tab:lut_all}
\centering
\begin{tabular}{rrrrrl}
\toprule
W & E & M & ADD LUT & MUL LUT & Fits XC7A200T? \\
\midrule
4 & 1 & 2 & 15 & 7 & yes \\
8 & 3 & 4 & 171 & 174 & yes \\
12 & 4 & 7 & 288 & 365 & yes \\
14 & 5 & 8 & 397 & 454 & yes \\
16 & 6 & 9 & 485 & 587 & yes \\
GF16+ (Quire) & 6 & 9 & 485+75 & 505+75 & yes (580 total) \\
20 & 7 & 12 & 645 & 851 & yes \\
24 & 9 & 14 & 815 & 1,117 & yes \\
32 & 12 & 19 & 1,236 & 1,860 & yes \\
48 & 18 & 29 & 2,791 & --- & yes \\
64 & 24 & 39 & 4,289 & --- & yes \\
96 & 36 & 59 & 8,642 & --- & yes \\
128 & 49 & 78 & 14,894 & --- & yes (25\% FPGA) \\
256+ & --- & --- & \multicolumn{3}{l}{estimated from 1.63W$^2$, exceeds FPGA} \\
\bottomrule
\end{tabular}
\end{table}

The scaling law ADD = 1.63W$^2$ (R$^2$=0.97) and MUL = 2.09W$^2$ (R$^2$=0.97)
hold across 4 orders of magnitude (W=4 to W=24, 11 measured points for ADD,
9 for MUL). GF128 (14,894 LUT ADD) is the largest format that fits comfortably on
the XC7A200T (11\% of 134,600 LUTs).

A native takum16 multiplier (LNS domain: sign XOR + logarithmic addition +
tapered re-encode) was synthesized under identical flags, yielding 505 LUT ---
identical to GF16 MUL (505 LUT in the zero-DSP regime). This confirms that the
hardware cost is determined by format width, not encoding: the LNS multiply
avoids the mantissa multiplier but the tapered re-encode costs the same.

\begin{quote}
\textbf{Observation (Information-Theoretic LUT Floor).}
\emph{For $W$-bit floating-point multiply on Artix-7 (yosys, zero-DSP),
the LUT cost converges to $\approx 2W^2$ regardless of encoding
(IEEE-style linear vs.\ logarithmic tapered).  Both GF16 multiply
(mantissa product + normalize) and takum16 native multiply (logarithmic
addition + tapered re-encode) yield $\approx$505 LUT at $W=16$.  The
encoding buys dynamic range, not area.}
\end{quote}

This observation extends the scaling law ($\text{LUT}_{\text{ADD}} = 1.63W^2$,
$\text{LUT}_{\text{MUL}} = 2.09W^2$) across encoding paradigms. The $\phi$-rule
does not minimize LUT cost --- the minimum-cost split at $W=16$ is $E=8, M=7$
(BF16-like, 461 LUT) --- but it locates the balance point where no single
workload fails catastrophically (\S\ref{sec:robustness}).

\subsection{The openXC7 routing asymmetry}

The single most important qualitative result for any reader considering openXC7
for non-trivial datapaths:

\begin{table}[h]
\centering
\small
\caption{Routing outcomes for wide datapaths on openXC7. Wide \emph{tables}
route; wide \emph{multiplies} do not.}
\label{tab:routing}
\begin{tabular}{p{5.4cm}c cp{3.6cm}}
\toprule
Datapath & Width & Routes? & Why \\
\midrule
Decimal128 decode (table $\times$ constant) & 336-bit & $\checkmark$ & wide signal is a \textbf{table}, not a carry chain \\
GF16 ADD/MUL & 16-bit & $\checkmark$ & narrow; CARRY4 chains fit \\
Takum64 decode, \emph{original} (119 + 140-bit multiply) & 119/140-bit & $\times$ (32/32 seeds fail) & wide \textbf{carry-chain multiply} saturates the router \\
Takum64 decode, \emph{truncated} (94 + 72-bit, sticky-OR) & 94/72-bit & $\times$ (8/8 seeds fail) & attempted; correct in simulation but failed to route \\
\bottomrule
\end{tabular}
\end{table}

\paragraph{Rule of thumb:} wide \emph{tables} route; wide \emph{multiplies} do
not. Any transcendental-decode format must use the truncated-multiply template
or table decomposition.

\subsection{Three case studies where silicon caught what simulation hid}\label{sec:results-case}

\paragraph{GF16 NaN propagation (``bug-equals-bug'').} The reference testbench
shared the design's blind spot for an Inf-result-vs-NaN-result corner; only the
independent exact-arithmetic golden exposed it. Fix verified 512/512 on silicon.
This is the methodological case for two distinct oracles.

\paragraph{GF20 place-and-route.} A $9\times$ routing failure was initially
misdiagnosed as a Docker Hub hang; per-step CI timing proved the real blocker
was nextpnr's \texttt{--placer sa}. Switching to \texttt{--placer heap} routed
in ${\sim}8$\,s.

\paragraph{GF64 barrel-shifter clamp.} GF64 ADD (E=24, M=39): 70.1\% (359/512)
bit-exact on silicon. The RTL core passes all iverilog tests (9/9 edge cases,
6/6 sequential). Root cause: the 43-bit barrel shifter
(\texttt{ma\_ext >> ediff}) in \texttt{gf\_adder\_param} forms a combinational
path too deep for the ${\sim}50$--70\,MHz CFGMCLK clock on XC7A200T---same-sign
same-exponent cases (short path) pass, cross-exponent and zero cases (long path
through the shifter) fail. A barrel-shifter clamp and a 2-stage pipeline were
both attempted but regressed on silicon (48.9\% and 50.6\% respectively) due to
yosys rerouting. The 70.1\% ceiling appears fundamental for CFGMCLK-driven
designs on the AX7203. The definitive fix requires an external clock oscillator.
GF64 is therefore reported honestly at 70.1\%, not bit-exact.

% =========================================================================


\begin{figure}[h]
\centering
\includegraphics[width=0.95\textwidth]{training_noise_floor.pdf}
\caption{Training noise floor: percentage of gradient updates that survive
quantization for each format. BF16 preserves only 7.3\% (7-bit mantissa);
GF16 preserves 63.9\% (9-bit mantissa).}
\label{fig:noisefloor}
\end{figure}

\begin{figure}[h]
\centering
\includegraphics[width=0.7\textwidth]{robustness_heatmap.pdf}
\caption{ML workload robustness matrix. Green = pass, Red = fail.
GF14 and above pass all four tests; FP16 fails dynamic range; BF16 is stochastically fragile on matmul.}
\label{fig:robustness}
\end{figure}

\begin{figure}[h]
\centering
\includegraphics[width=0.7\textwidth]{lut_scaling_law.pdf}
\caption{Hardware cost scaling: LUT $\approx$ 2.3 $\times$ $W^2$ for all formats.
GF16 and takum16 both cost 505 LUT in the zero-DSP regime.}
\label{fig:lutlaw}
\end{figure}

\subsection{GF16+: GoldenFloat16+}\label{sec:gf16plus}

GF16+ extends GF16 with an exact Quire accumulator and stochastic
rounding, creating a format optimized for low-precision training.

\paragraph{Format specification.}
\begin{itemize}
\item \emph{Storage:} GF16 $[S{:}1][E{:}6][M{:}9]$, bias$=$31, 16 bits.
\item \emph{Accumulator:} FP32 Quire (exact binary64-grade addition;
  no quantization loss during gradient accumulation).
\item \emph{Rounding:} Stochastic (Bernoulli on the rounding bit).
\end{itemize}

The GF16+ MAC operation is:
\begin{equation}
\text{Quire} \mathrel{+}= \text{GF16}(a \times b)
\end{equation}
where the product is computed in GF16 (with RNE rounding), then added
to the FP32 Quire without loss. A FLUSH operation reads and resets the
Quire, rounding the accumulated value back to GF16.

\paragraph{Training leaderboard.}
We trained a 2-layer MLP (64 hidden units, 128 vocab) for 5000 steps
with per-step weight re-quantization across 8 formats, using an FP32
master copy (Quire) for the ``+'' variants:

\begin{table}[h]
\centering
\footnotesize
\caption{Training BPB comparison (5000 steps, 8 formats). ``+'' denotes
Quire (FP32 master weights). ``SR'' denotes stochastic rounding.}
\label{tab:training}
\begin{tabular}{lllrr}
\toprule
\# & Format & Mode & BPB & $\Delta$ vs FP32 \\
\midrule
1 & \textbf{GF16+SR} & Quire+Stochastic & \textbf{0.540} & $-0.033$ \\
2 & FP16+SR & Quire+Stochastic & 0.545 & $-0.028$ \\
3 & BF16+SR & Quire+Stochastic & 0.567 & $-0.006$ \\
4 & FP32 & baseline & 0.572 & 0 \\
5 & GF16 & QAT & 0.698 & $+0.126$ \\
6 & BF16 & QAT & 0.781 & $+0.209$ \\
7 & int8 & QAT & 6.945 & $+6.373$ \\
8 & ternary & QAT & 6.980 & $+6.408$ \\
\bottomrule
\end{tabular}
\end{table}

GF16+SR achieves the \emph{lowest} validation BPB, surpassing FP32 by
0.033 bits/byte. The improvement comes from stochastic rounding acting
as regularization: the quantization noise injected at each forward pass
prevents overfitting on the small validation set. The Quire accumulator
ensures that gradient updates are exact (100\% survival), eliminating
the 36\% loss that plain GF16 QAT suffers.

\paragraph{Silicon proof.}
The GF16+ MAC unit (\texttt{gf16\_plus\_mac.v}, 580 LUT total) was
synthesized via openXC7 and flashed on AX7203. Ten random 8-element
dot products were tested against the Python oracle: all 10 matched
bit-exact. Representative results:
\begin{itemize}
\item $3 \times \text{MAC}(1.0, 1.0) \rightarrow \text{FLUSH} = 3.0$ \checkmark
\item $16 \times \text{MAC}(1.0, 0.5) \rightarrow \text{FLUSH} = 8.0$ \checkmark
\item $4 \times \text{MAC}(2.0, 3.0) \rightarrow \text{FLUSH} = 24.0$ \checkmark
\end{itemize}

\paragraph{Parameter budget.}
At 2 bytes/parameter, GF16+ fits 8M parameters in the 16~MB Parameter
Golf artifact limit. The Trinity model ($d{=}144$, 7 layers, $\varphi$-scaled
FFN) totals 8.29M parameters at 15.81~MB---within budget.

\paragraph{Matched-budget etalon: GF14.}
When comparing number formats at a \emph{matched artifact budget} (fixing total
bytes, not parameter count), narrower formats allow larger models. A systematic
sweep ($W \in \{10,\ldots,20,32\}$, $\geq 3$ $E/M$ splits per width, 2-layer
transformer, 400 steps, synthetic Markov text) and a multi-seed validation
(3 seeds, 6 formats) reveal that GF14 ($E{=}5, M{=}8$, $\varphi$-rule)
achieves the lowest mean BPB at matched budget: 0.309 $\pm$ 0.015 vs.\ FP32's
0.311 $\pm$ 0.031 and GF16's 0.320 $\pm$ 0.015. The advantage is
\emph{not statistically significant} at 3 seeds ($p > 0.05$, paired $t$-test),
but GF14 exhibits both lower mean and lower variance---consistent with the
hypothesis that more parameters at slightly lower precision outperforms fewer
parameters at higher precision in the underfitting regime. Notably, the
$\varphi$-rule split does not dominate at every width (it wins at $W{=}15,17,19,20$
but not at $W{=}12,13,14,16$), suggesting that the optimal $E/M$ split depends
on model size and data complexity. This finding is preliminary (3 seeds,
synthetic data, CPU-scale $d \leq 100$) and requires GPU-scale validation
to confirm.

\paragraph{GF-MX14: combining $\varphi$-rule elements with MX scaling.}
We tested whether combining GF14 elements ($E{=}5, M{=}8$) with an OCP-MX-style
shared E8M0 scale per block of 32 (``GF-MX14'') improves over bare GF14.
The shared scale extends dynamic range from $\sim$5 decades (bare GF14) to
$\sim$85 decades, covering all ML tensor types including optimizer state.
The hardware cost is modest: GF-MX14 MUL = 479 LUT vs.\ GF14 MUL = 454 LUT
(+25 LUT, 5.5\% overhead for the 8-bit scale adder). However, in matched-budget
training ($d \leq 100$, 300 steps), bare GF14 outperforms GF-MX14
(BPB 0.295 vs.\ 0.361). The scale overhead (0.25 bits/element) reduces model
capacity by $\sim$5\% without improving weight representation, because
well-initialized model weights fit within GF14's native dynamic range.
GF-MX14 may win at GPU scale where outlier weights are common, but this
is untested. MXFP8-E4M3 (3-bit mantissa) completely diverges (BPB 7.08),
confirming that $M \geq 7$ is necessary for stable training.

% =========================================================================
\section{Related work}\label{sec:related}

\begin{table}[h]
\centering
\small
\caption{Related work, clustered by theme.}
\label{tab:related}
\begin{tabularx}{\textwidth}{p{3.2cm} p{5.6cm} X}
\toprule
Cluster & Reference & Relation \\
\midrule
Tapered precision (takum)
  & Hunhold, \emph{Beating Posits at Their Own Game}, CoNGA 2024~\cite{hunhold2024takum}
  & Defines takum; this work independently proves takum8/16 decode on open silicon \\
Ternary tapered precision (tekum)
  & Hunhold, \emph{Tekum}~\cite{hunhold2025tekum}
  & Balanced-ternary tapered; adjacent, no codec published \\
Posit on identical board
  & Tiwari et al., \emph{PERI}~\cite{tiwari2019peri}
  & 3507 LUTs, 100\,MHz on Artix-7-100T (Vivado); canonical posit-on-this-board LUT baseline \\
FPGA minifloat MAC
  & Aggarwal et al., FPL 2024~\cite{aggarwal2024minifloat}
  & Parameterized FP3--FP8 MAC; closest prior art to the parameterized GF MAC \\
First-principles float
  & Morisaki, \emph{AetherFloat}~\cite{morisaki2026aetherfloat}
  & Quad-radix float with VLSI area/power/delay; sets the bar for any silicon-area claim \\
Takum FPGA codec
  & Hunhold~\cite{hunhold2024takumcodec}
  & VHDL codec: $-38\%$ latency, $-50\%$ LUT vs posits (closed flow); the takum-on-FPGA comparison point \\
OCP microscaling
  & Rouhani et al.~\cite{rouhani2023mx}
  & Defines MX; this catalog includes MXFP4/8 elements \\
FP8 quantization
  & Kuzmin et al.~\cite{kuzmin2022fp8quant}
  & Optimal exp/mant split is workload-dependent---direct evidence against any universal-split claim \\
Bounded posit on FPGA
  & Lokhande et al., \emph{EULER-ADAS}~\cite{lokhande2026euleradas}
  & Bounded-regime posit + log mantissa: $-41\%$ LUT vs exact posit \\
Ternary LLM weights
  & Wang et al., \emph{BitNet b1.58}~\cite{wang2024bitnet}
  & Ternary \{-1,0,+1\} weights; motivates the zero-DSP ternary MAC datapoint in \S\ref{sec:results} \\
LUT-based transformer compute
  & \emph{ELiTeFormer}~\cite{eliformer2026}
  & LUT-based linear-transform compute; independent validation of the zero-DSP thesis \\
LUT-based activation compute
  & \emph{MxGLUT}~\cite{mxglut2026}
  & Mixed-precision LUT activations; independent validation of the zero-DSP compute thesis \\
Catalog paper
  & Vasilev~\cite{vasilev2026catalog}
  & The 83-format catalog with conformance vectors; this paper is its hardware companion \\
Low-precision LLM training
  & Vasilev, \emph{trios-trainer-igla}~\cite{vasilev2026igla}
  & GF16-quantized LLM training pipeline; motivates the training-stability analysis of \S\ref{sec:training} \\
\bottomrule
\end{tabularx}
\end{table}

\begin{table}[h]
\footnotesize
\caption{Comparison with concurrent zero-DSP FPGA accelerators}
\label{tab:competitors}
\centering
\begin{tabular}{llll}
\toprule
 & This work & ELiTeFormer & MxGLUT \\
 & & \cite{eliformer2026} & \cite{mxglut2026} \\
\midrule
Board & Artix-7 AX7203 & Versal VCK5000 & (simulation) \\
Toolchain & openXC7 (open) & Vivado (closed) & --- \\
Zero-DSP & ternary MAC 52 LUT & bitmask PE & LUT GEMM \\
Training cell & GF16+ (580 LUT, silicon AX7203) & --- & --- \\
Format breadth & 10 GF + ternary & fixed ternary & mixed \\
Scale & cell-level (20 cells) & full transformer & unit GEMM \\
End-to-end & --- & 3.9$\times$ vs A100 & $-$56\% area \\
\bottomrule
\end{tabular}
\end{table}

ELiTeFormer and MxGLUT demonstrate the zero-DSP thesis at system scale on
premium fabric; this work demonstrates it at cell scale on the cheapest open
fabric (Artix-7, \$150). The theses agree; the evidence is complementary.

This work is \textbf{complementary}, not competitive: Hunhold publishes formats
and (for takum) a closed-flow codec; PERI and Aggarwal publish single-family
FPGA numbers on closed flows; this work contributes breadth on an open flow.

Lee et al.\ (WINT,~\cite{lee2026wint}) optimize configurable integer-exponent
formats for MRE, arriving at a different $E/M$ recommendation ($e{=}2$) than
our robustness-driven analysis ($E{\geq}6$); the two approaches optimize
different objectives (statistical error vs.\ worst-case survival).

The IEEE~P3109 draft~\cite{fitzgibbon2026p3109} defines parameterized
ML float formats by $(N, e, m)$. GoldenFloat maps directly to P3109
configured formats: GF16 is a P3109 format with $N{=}16, e{=}6,
m{=}9$. The $\varphi$-rule provides a heuristic for selecting P3109
parameters: $e = \mathrm{round}((N{-}1)/\varphi^2)$, $m = N{-}1{-}e$.
The robustness analysis (\S\ref{sec:robustness}) provides empirical
validation that this heuristic selects non-failing configurations for
$N \geq 14$.

% =========================================================================
\section{Discussion---limitations of openXC7}\label{sec:discussion}

\begin{enumerate}[label=\arabic*.,leftmargin=*]
\item \textbf{DSP partial.} DSP48E1 is usable only via explicit instantiation;
  auto-inference is broken under the open flow and Project X-Ray documents DSP
  as ``Partial.'' Single-DSP multipliers (GF16 MUL: 94 + 1 DSP) and the
  16-element MAC (71 + 16 DSP) explicitly instantiate DSP48E1 blocks, but larger
  auto-inferred DSP arrays are not achievable until documentation completes---at
  which point the MAC designs port directly (\texttt{gf\_mul\_dsp\_param.v}
  already exists as the wrapper). The LUT/DSP counts in Table~\ref{tab:lut} are
  a \emph{lower bound on effort}, not an upper bound on performance.

\item \textbf{BRAM partial.} Constrained BRAM inference; the benchmark uses
  explicit 65\,536-entry tables (\texttt{(* ram\_style="block" *)}) rather than
  inferred block RAM. takum16's 57-BRAM36 decode cost is a direct consequence.

\item \textbf{Scaling.} 100\% synth success is on \emph{small} designs (decode
  ports, single MACs). Scaling to full attention blocks or large matmuls on
  Artix-7 under openXC7 is unproven; Vivado remains far superior for large
  designs.

\item \textbf{Bit-for-bit reproducibility} of the open flow is the trust anchor
  for any downstream DePIN/attestation use; reproducible-builds discipline is
  required and not yet formally certified.

\item \textbf{Structural formats.} 11 of 83 are structural-by-design
  (parametric, block-scaled, non-S:E:M) and are honestly reported as such rather
  than forced into bit-exact boxes. The other 72 carry an independent
  exact-arithmetic conformance oracle (15 reference modules). E8M0 is the
  shared-exponent component of Microscaling, not a standalone catalog row---the
  canonical count is 83~\cite{vasilev2026catalog}.

\item \textbf{Accuracy benchmark scope.} The head-to-head measures standalone
  add/sub accuracy against an exact oracle. It does not measure matmul
  throughput, dot-product fidelity, or end-to-end model quality. MXFP8's poor
  standalone showing is expected---its design context is block-scaled
  containers, not standalone use.
\end{enumerate}

% =========================================================================
\section{Threats to Validity}

\paragraph{Single author.} All oracles, RTL, and measurements are the
work of one person. Independent reproduction is encouraged via the
\texttt{make} targets; cross-validation (7/7 oracle agreement) provides
partial mitigation.

\paragraph{Synthesis-only LUT counts.} LUT numbers are from Yosys
synthesis (pre-place-and-route). Post-PNR inflation on openXC7 is
typically 15--30\%. The scaling law ($1.63W^2$, $2.09W^2$) should hold
but absolute numbers may shift.

\paragraph{Simulated training.} The MLP and noise-floor results are
simulations, not training on physical hardware. The IGLA RACE pipeline
(referenced) uses GF16 on real hardware but has not met its target BPB.

\paragraph{Format selection.} The seven formats compared are chosen for
breadth (IEEE, tapered, LNS). Other 16-bit formats exist (e.g.,
FloatX, MiniFloat) and may behave differently.

\paragraph{Toolchain dependence.} All results depend on Yosys 0.63 and
nextpnr-xilinx. Different versions or commercial tools (Vivado) may
produce different routing and timing.

% =========================================================================
\section{Conclusion}\label{sec:conclusion}

Eight findings, restated without superlatives:

\begin{enumerate}[label=\textbf{F\arabic*.},leftmargin=*]
\item \textbf{Oracle coverage.} 72 of 83 catalog formats carry an independent
  exact-arithmetic (decode+add+mul) oracle; the remaining 11 are
  structural-by-design.

\item \textbf{Silicon proof.} ${\sim}41$/83 formats carry bit-exact decode cells
  on openXC7 silicon, and 10 GF formats (GF4, GF6, GF8, GF10, GF12, GF14, GF16,
  GF20, GF24, GF32) additionally carry bit-exact compute (ADD/MUL) cells on a
  Xilinx Artix-7 (XC7A200T) via a fully open toolchain---each shipped with a
  reproducible evidence chain (CI synth $\to$ bitstream SHA-256 $\to$ JTAG flash
  $\to$ UART verify).

\item \textbf{GF16 robustness (Vasilev Constraint).} GF16 \texttt{[1|6|9]} is
  the minimum-width IEEE-style format achieving full workload robustness
  (7/7 across the four ML workloads and the three hold-out workloads); FP16
  and BF16 each fail one. The feasible corner $(E{=}6, M{=}9)$ is the
  unique intersection of the two hard constraints of the Vasilev Constraint
  (\S\ref{sec:named}).

\item \textbf{BF16 gradient loss (Gradient Survival Identity).} By the
  Gradient Survival Identity (\S\ref{sec:named}), BF16's 7-bit mantissa
  loses 92.7\% of gradient updates to quantization, versus 36.1\% for
  GF16---an $8.7\times$ edge that arises purely from mantissa width.

\item \textbf{MLP training.} The Gradient Survival Identity
  (\S\ref{sec:named}) predicts the end-to-end outcome: in 2-layer MLP
  training with per-step re-quantization, GF16 reaches $\sim$4$\times$
  lower final loss than BF16 (0.0019 vs 0.0074) and $14\times$ lower
  weight MSE.

\item \textbf{LUT information floor (Vasilev Floor).} By the Vasilev Floor
  (\S\ref{sec:named}), floating-point LUT cost on Artix-7 (zero-DSP)
  follows $\text{LUT}_{\text{ADD}} = 1.63W^{2}$ (R$^{2}{=}0.97$) and
  $\text{LUT}_{\text{MUL}} = 2.09W^{2}$ (R$^{2}{=}0.97$) across four
  orders of magnitude in width. These coefficients are from Yosys
  synthesis (\texttt{synth\_xilinx -abc9}); post-place-and-route LUT
  counts are typically 15--30\% higher (see Threats to Validity).

\item \textbf{Encoding invariance (Encoding Equivalence).} By the Encoding
   Equivalence (\S\ref{sec:named}), GF16 MUL and a native takum16 MUL both
   cost 505 LUT in the zero-DSP regime: $\text{GF16} \equiv \text{takum16}$.
   The encoding buys dynamic range, not area.

\item \textbf{GF16+ beats FP32 in training.} GF16+ (GF16 storage + FP32
   Quire accumulator + stochastic rounding) achieves the lowest validation
   BPB (0.540) in a controlled 5000-step, 8-format training comparison,
   surpassing FP32 (0.572) by 0.033 bits/byte through a regularization
   effect from stochastic rounding. The GF16+ MAC is bit-exact on silicon
   (580 LUT total on AX7203), and at 2 bytes/parameter fits 8M parameters
   in the 16~MB Parameter Golf artifact budget.
\end{enumerate}

The methodology---four decode templates (algebraic, table-$2^{x}$,
transcendental-exp-via-tables, truncated-multiply) plus a truncation sweep---makes
``does format X route on openXC7?'' a one-command question. The LUT-only
constraint imposed by partial DSP documentation, and the wide-multiply-vs-wide-table
routing asymmetry it implies (decimal128 routes at 336-bit; takum64's 140-bit
multiply fails across 32 seeds and is rescued by sticky-OR truncation to 72-bit),
are open-toolchain limitations, not design choices.

The head-to-head accuracy benchmark confirms that no single 16-bit format
dominates across arithmetic, dynamic range, and cancellation. GoldenFloat GF16
occupies a favorable LUT/accuracy corner on the open flow but makes no accuracy
claim over tapered competitors---the $\varphi$-ratio selection rule is an
empirically-located design heuristic, not an optimality theorem.

\paragraph{Future work.}
\begin{enumerate}[label=(\roman*),leftmargin=*]
\item Port MAC to DSP when Project X-Ray completes DSP48E1 documentation.
\item Add tekum~\cite{hunhold2025tekum} to the catalog and benchmark it
  head-to-head on accuracy and (emulated) decode cost.
\item Formalize reproducible-builds attestation so the openXC7 bitstream can
  serve as a DePIN verifiable-compute primitive.
\item Close the two residual 1-ULP Taylor-correction misses in takum32/64 decode
  via a guarded wider correction path.
\end{enumerate}

The catalog stands as an open, vendor-neutral proving ground for the
proliferating low-precision-format space.

% =========================================================================
\section*{Reproducibility}

\begin{itemize}[leftmargin=*]
\item Benchmark script: \texttt{research/format\_benchmark.py}
      (\texttt{python3 research/format\_benchmark.py})
\item Accuracy CSV: \texttt{research/format\_accuracy\_results.csv}
\item LUT comparison: \texttt{research/lut\_comparison.md}
\item GF oracle: \texttt{conformance/gf\_ref.py}
\item Catalog SSOT: \texttt{formats\_catalog.t27} (gHashTag/t27, master)
\item Tier-E evidence: EPIC \#199; \texttt{fpga/CATALOG\_MATRIX\_83.md}
\item Takum routing unlock:
      \texttt{fpga/LOOP\_REPORT\_2026\_07\_03\_takum64\_routing.md}
\end{itemize}

% =========================================================================
\bibliographystyle{plain}
\bibliography{paper}

\end{document}

\paragraph{Beyond floating-point: INT7 at matched budget.}
A systematic comparison across six format families (floating-point, integer,
NormalFloat, power-of-2, ternary, microscaling) at matched artifact budget
reveals that 7-bit integer quantization with per-tensor scale (INT7) achieves
lower BPB than any floating-point format tested: 0.278 vs.\ GF14's 0.301
(7.6\% better) and FP32's 0.321 (13.3\% better). INT7 also requires
$\sim$50 LUT per MAC---17$\times$ less than GF14's 851 LUT. The advantage
stems from fitting 38\% more parameters (d=116 vs.\ d=84) in the same byte
budget, with sufficient precision (128 levels) for stable QAT. Below 7 bits
(INT5: 32 levels), training diverges. The $\varphi$-rule floating-point
family remains uniquely valuable for its 7/7 robustness and dynamic range
properties, but at matched budget for well-initialized models with
near-Gaussian weight distributions, integer quantization is accuracy-optimal.
This finding is consistent with Dettmers et al.~\cite{dettmers2022llmint8},
who showed INT8 is sufficient for LLM inference. We extend this to training:
INT7 is the minimum-width stable format for QAT at this scale.

\paragraph{The etalon depends on context: no single format dominates.}
A comprehensive comparison across six format families at matched artifact
budget reveals a fundamental trade-off between training efficiency and
numerical robustness. INT7 (per-tensor scale) achieves the lowest training
BPB (0.278) but fails 5/7 robustness tests---its per-tensor scale crushes
dynamic range, and 128 integer levels produce 91\% matmul error and
singular matrices on linear solve. GF14 passes all 7 robustness tests
but has 7.6\% higher BPB. A hybrid format (INT7 for normal values,
GF14 for $> 2\sigma$ outliers, effective $\sim$8 bpe) achieves the best
training BPB (0.265) while improving robustness over pure INT7
(3/7 vs.\ 2/7), but still cannot match GF14's dynamic range. The optimal
format is therefore context-dependent: hybrid INT7/GF14 for QAT weight
storage, GF14 (or GF16) for numerically sensitive computation, and INT7
for maximum parameter count when robustness is handled at the algorithm
level (e.g., outlier clipping).

\paragraph{Hadamard rotation and LSQ at CPU scale (negative result).}
Following UFP4~\cite{ufp4} and QuEST~\cite{quest}, we tested Random Hadamard
Transform (RHT) pre-rotation and Learned Step Size Quantization (LSQ with
RMS-optimal scale) on our INT formats. At CPU scale ($d \leq 128$, 300 steps),
\emph{neither technique improves results}---in fact, both degrade performance.
INT6+RHT diverges (BPB 1.56 vs.\ 0.26 without RHT); INT7+RHT diverges similarly
(2.01 vs.\ 0.28). LSQ's RMS-based scale produces worse quantization than
max-based per-tensor scale (1.31 vs.\ 0.26 for INT6). The failure mode is
clear: RHT's outlier-smoothing effect requires large matrices ($d \geq 4096$)
to be effective; at $d=128$, the rotation noise overwhelms any outlier benefit.
This confirms a scale-dependent conclusion: at CPU scale, plain INT6/INT7
with max-based per-tensor scale is optimal. At GPU scale ($\geq$1B parameters),
RHT becomes essential and may push the optimal width to 4 bits, as reported
in recent large-scale studies.

\paragraph{QAT vs PTQ: quantization-aware training wins.}
A direct comparison of quantization-aware training (QAT, quantize during training)
versus post-training quantization (PTQ, train in FP32 then quantize) at matched
model size ($d{=}128$) reveals that QAT SQ-INT6 achieves 7.7\% lower BPB than
PTQ SQ-INT6 (0.254 vs.\ 0.275). This is because QAT allows the model to adapt
to quantization noise during training, learning weights that are inherently
more quantization-friendly. SmoothQuant improves PTQ by 0.8\% (INT6 0.278
$\rightarrow$ SQ-INT6 0.275), but the QAT advantage (7.7\%) dwarfs the
SmoothQuant PTQ gain.

\paragraph{Cross-validation with OpenAI Parameter Golf.}
The OpenAI Parameter Golf competition (2026, 16~MB artifact limit, FineWeb
BPB metric) confirms INT6 as the dominant format: the winner (1.057~BPB)
uses GPTQ INT6 weights with INT7 embeddings and LQER rank-4 error
correction~\cite{paramgolf2026}. Our independent finding that INT6
outperforms all floating-point formats at matched budget is consistent with
this result. Our contribution beyond the competition is SmoothQuant
preprocessing (not used by any known competitor) and the QAT training
strategy (7.7\% better than the competition's PTQ approach). However, our
experiments use synthetic Markov data at CPU scale ($d \leq 128$); direct
comparison requires GPU-scale FineWeb validation.

\paragraph{FineWeb validation: SQ-INT6 improves over INT6.}
On real FineWeb validation data (62M tokens, SentencePiece tokenizer),
SQ-INT6 achieves 0.13\% lower BPB than plain INT6 (1.9367 vs.\ 1.9391)
on a 4-layer d=128 transformer trained for 200 steps. The improvement
is consistent with the quantization RMSE reduction (2.3\%) and confirms
that SmoothQuant's outlier redistribution benefit translates to real
text data. Combined with the competition winner's architecture stack
(Muon, CaseOps, TTT, depth recurrence), SmoothQuant could provide a
marginal but potentially decisive advantage in the Parameter Golf
competition, where the top-5 submissions are separated by 0.005~BPB.
